\section{The theta algebra and its projective family}\label{theta:section}

The scalar graphs make a computable part of the directed category.  Their
products will provide the ring over which the new brane's cohomology becomes
a graded module.  We first compute these products over the formal coefficient
ring.  We then realize the ring, by an abstract field embedding, as a classical
theta section ring.  This second step uses convergence only of the explicit
theta series.

Keep the symmetric complex matrix $M$ commuting with $D$, the integer
$h_0\geq3$, and $R=\kappa[[t]]$ from Lemma~\ref{cat:residue}.  Thus
$\Lambda\hookrightarrow\kappa$, the field $\kappa$ is algebraically closed,
and $|\kappa|\leq|\C|$.  Throughout this section $G_i$ means the layer-zero
graph
\[
 p=h_0i q+\nabla f_i(q),\qquad i\geq1.
\]
The periodic potentials have sufficiently small Hessians.  We use the
integer grading on arrows between these graphs, before folding modulo four.

\subsection{Intersections and normalized triangle weights}

A lifted sheet of $G_i$ has equation
$p=h_0i q+\nabla f_i(q)+a_i$, with $a_i\in\Z^8$.  If $i>j$, put
$m=h_0(i-j)$ and $u=(a_j-a_i)/m$.  Its intersection with a lifted sheet of
$G_j$ is the unique solution of
\begin{equation}\label{theta:intersection}
 m(q-u)+\nabla(f_i-f_j)(q)=0.
\end{equation}
Indeed the left side is the gradient of the proper strictly convex function
$m|q-u|^2/2+f_i(q)-f_j(q)$.  The intersection generators are consequently
indexed by
\[
 \bar u\in (m^{-1}\Z^8)/\Z^8,
\]
and there are $m^8$ of them.  All have integer degree zero.  A holomorphic
strip between these two graphs lifts to fixed sheets: its ends are their
same unique intersection.  Each sheet is exact on the cover, so Stokes'
theorem makes the strip's area zero.  Thus every such strip is constant.
In particular the arrow differential vanishes; the grading also makes
every operation on decreasing-slope words except the binary product zero.
The binary product is therefore associative on the actual arrow spaces.

For $i>j>k$, set
\[
 m=h_0(i-j),\qquad m'=h_0(j-k),\qquad
 u=\frac{a_j-a_i}{m},\quad v=\frac{a_k-a_j}{m'},\quad
 w=\frac{mu+m'v}{m+m'}.
\]
For a symmetric matrix $U$ define
\begin{align}
 T^U_{m,m'}(u,v)
 &:=\frac12\bigl(mu^tUu+m'v^tUv-(m+m')w^tUw\bigr)\notag\\
 &=\frac{mm'}{2(m+m')}(u-v)^tU(u-v).
 \label{theta:quadratic}
\end{align}
These quadratic expressions are precisely the triangle weights after a
change of basis.  Here is the change, including the boundary primitive
terms.

For $B=dq^tM\wedge dp$, a primitive on the cover is $-p^tM\,dq$.
On the graph use the prescribed relative primitive
$\xi_i=-(\nabla f_i)^tM\,dq$.  Hence the potential of
$-p^tM\,dq-\xi_i$ on sheet $i$ is
\[
 -\frac{h_0i}{2}q^tMq-a_i^tMq.
\]
For $\omega$ the analogous potential is
$-h_0i|q|^2/2-a_i^tq-f_i(q)$.  At the intersection $q_{ij}$ the difference
of the sheet-$i$ and sheet-$j$ potentials is respectively
$mu^tMu/2+\delta^B_{ij}$ and $m|u|^2/2+\delta^\omega_{ij}$, where
\begin{equation}\label{theta:corner-corrections}
 \begin{split}
 \delta^B_{ij}&=-\frac m2(q_{ij}-u)^tM(q_{ij}-u),\\
 \delta^\omega_{ij}&=-\frac m2|q_{ij}-u|^2-f_i(q_{ij})+f_j(q_{ij}).
 \end{split}
\end{equation}
These quantities are unchanged by simultaneous integral translation of
$q_{ij}$ and $u$, and thus depend only on the projected corner.  In the
counterclockwise boundary integral the two input differences occur with
positive sign and the output difference with negative sign.  Replacing
each intersection generator by
\begin{equation}\label{theta:normalization}
 \qnov^{-\delta^\omega_{ij}}
       \exp(-t\delta^B_{ij})\, [q_{ij}]
\end{equation}
cancels all three corrections.  A triangle with these lifted boundaries
therefore contributes its sign times
\[
 \qnov^{T^{1_8}_{m,m'}(u,v)}
       \exp\bigl(tT^M_{m,m'}(u,v)\bigr).
\]
The factors in Equation~\eqref{theta:normalization} are units of $R$.
Their Novikov valuations are fixed finite numbers; the coefficient-limit
construction is unchanged by this normalization.

\begin{lemma}[The count of a lifted triangle]\label{theta:triangle-count}
For every fixed triple of lifted sheets with $i>j>k$, the unweighted
oriented triangle count has absolute value one.  After choosing the
intersection orientation bases equivariantly, its sign depends only on
$(i,j,k)$.
\end{lemma}

\begin{proof}
Shrink the potentials $f_i,f_j,f_k$ to zero.  Their pair differences remain
strictly convex, so all lifted corners continue uniquely.  Deform the
triangle data to the standard complex structure $J_*$ at the affine end.
The primitive computation above bounds area in this fixed lifted class
uniformly along the deformation.  There are no spheres or disks on a
single graph, and no nonconstant strips can break off, by the uniqueness
and exactness just proved.  A three-marked disk has no domain boundary
without such bubbling or breaking.  Thus the parameterized count is a
compact oriented cobordism.  If necessary, achieve transversality by
Hamiltonian one-form variations supported in the triangle interior,
fixed at the regular endpoints and at the strip ends.  The usual
continuation energy estimate adds only a uniformly bounded curvature
term; these variations introduce none of the excluded bubbles.  This
also handles parameter values with triple constant maps.  Anchoring one
corner fixes compatible lifts throughout the compactness argument.

At the affine endpoint the three vertices have real parts $w,u,v$ in
boundary order.  Their differences are complex scalar multiples of the
real vector $u-v$.  If $u\neq v$, they bound a positively oriented straight
triangle in a complex affine line, with area
$T^{1_8}_{m,m'}(u,v)>0$.  Uniformization gives a holomorphic triangle with
the three prescribed boundary marks.  If $u=v$, the three vertices
coincide and the solution is constant.

Any two solutions differ by a holomorphic map with boundary on the three
linear tangent planes and zero corner limits.  The planes have their
quadratic exact primitives, and Stokes' theorem gives this difference
zero area.  It is therefore zero.  The same argument gives zero kernel
for the linearized operator with the three domain marks fixed.  Its
index is zero by the degree calculation, so it is surjective.  This
proves regularity, including for the constant solution, and the asserted
absolute count.  At the affine endpoint the tangent planes, framed spin
data, and linearized operators depend only on the slopes.  Choose the cap
bases using these common frames, and transport them along the deformation.
Deck translation preserves these choices.  The resulting sign is thus
independent of the offsets.
\end{proof}

Write this sign as $\epsilon(i,j,k)\in\{1,-1\}$.  The unsigned formulas
are associative: for three inputs, the two sums range over the same
three lifted points modulo common translation, and their exponents agree
by telescoping the quadratic expression in
Equation~\eqref{theta:quadratic}.  The sums can be rearranged because the
quadratic area gives only finitely many terms below each energy bound.
Their iterated products are nonzero at $t=0$, where every coefficient is
a sum of positive Novikov monomials.  Comparing this associativity with
the signed Floer associativity gives
\[
 \epsilon(i,j,k)\epsilon(i,k,l)
   =\epsilon(j,k,l)\epsilon(i,j,l)\qquad(i>j>k>l).
\]
Rescale every arrow $G_i\to G_j$ with $j>1$ by
$\epsilon(i,j,1)$, and leave arrows to $G_1$ unchanged.  The displayed
identity makes every product sign $+1$.  Both the monomial normalization
and this sign normalization commute with reduction modulo $t$.

\subsection{The ring and its cohomology modules}

Define the graded $R$-module $\mathscr S$ by $\mathscr S_0=R$ and
\[
 \mathscr S_\ell=
 \bigoplus_{\bar u\in ((h_0\ell)^{-1}\Z^8)/\Z^8}
       R e_{\ell,\bar u}\qquad(\ell\geq1).
\]
Let $1\in R$ be its multiplicative identity.  For
$m=h_0\ell$ and $m'=h_0\ell'$, define
\begin{equation}\label{theta:product}
 e_{\ell,\bar u}e_{\ell',\bar v}
 =\sum_{\substack{u\in\bar u+\Z^8,\ v\in\bar v+\Z^8\\
                         \text{modulo simultaneous }\Z^8}}
  \qnov^{T^{1_8}_{m,m'}(u,v)}
  \exp\bigl(tT^M_{m,m'}(u,v)\bigr)
  e_{\ell+\ell',\overline{(mu+m'v)/(m+m')}}.
\end{equation}
In this formula an input residue denotes any representative of its coset;
the indicated quotient makes the sum independent of those representatives.
The output belongs to $((m+m')^{-1}\Z^8)/\Z^8$.
Fixing one representative of $u$ leaves a discrete lattice of possible
$v-u$.  The expression $T^{1_8}_{m,m'}$ is a positive definite quadratic
form on that lattice.  Consequently each coefficient of $t^a$ in each
output coefficient is a locally finite Novikov sum: inserting
$(T^M)^a/a!$ changes its coefficients, not its support.  Thus
Equation~\eqref{theta:product} is defined in $\Lambda[[t]]$ before applying
$\Lambda[[t]]\hookrightarrow R$.

\begin{proposition}[The directed theta ring]\label{theta:directed-ring}
Equation~\eqref{theta:product} makes $\mathscr S$ a commutative associative
graded $R$-algebra.  The normalized arrows give isomorphisms
\[
 \mathscr S_{i-j}\ \simeq\
       \Hom_{\mathcal D_R}(G_i,G_j),\qquad i>j,
\]
compatible with composition, and the arrow spaces have zero differential.
Each $\mathscr S_\ell$ is free of rank $(h_0\ell)^8$ for $\ell>0$.
All these identifications reduce to the unweighted ones at $t=0$.
\end{proposition}

\begin{proof}
Symmetry under exchanging $(m,u)$ and $(m',v)$ proves commutativity.
Associativity follows from the quadratic telescoping and convergence
argument above.  A compatible lift triple of sheets is determined, modulo
deck translation and common integral $p$-translation, by the two input
lifts in Equation~\eqref{theta:product}; its third corner is their weighted
average $w$.  Lemma~\ref{theta:triangle-count}, the action computation, and
the removal of $\epsilon$ therefore identify this multiplication with the
binary Floer operation.  The remaining assertions were established above.
\end{proof}

Set $\mathscr S_{\mathrm s}=\mathscr S/t\mathscr S$.
\begin{proposition}[The graded cohomology modules]\label{theta:modules}
For $r\in\Z/4$, the groups
\begin{equation}\label{theta:module-definition}
 V^r=\bigoplus_{j\geq1}H^r(\mathcal M_R(G_j)),\qquad
 V_{\mathrm s}^r=\bigoplus_{j\geq1}H^r(\mathcal M_0(G_j)),
\end{equation}
with the summands placed in ordinary degree $j$, are graded modules over
$\mathscr S$ and $\mathscr S_{\mathrm s}$, respectively.  Each ordinary
graded piece of $V^r$ is finite over $R$, and each such piece of
$V_{\mathrm s}^r$ is finite-dimensional over $\kappa$.  The actions are
induced by chain maps compatible with reduction modulo $t$.
\end{proposition}

\begin{proof}
The arrow $G_{j+\ell}\to G_j$ gives the right-module action from
Proposition~\ref{cat:module},
$\mathcal M_R(G_j)\to\mathcal M_R(G_{j+\ell})$.  It has degree zero, hence
is a chain map without an additional cohomological sign.  The module
identity makes the two binary actions for a pair of arrows chain
homotopic; passing to cohomology and using
Proposition~\ref{theta:directed-ring} gives the associative graded action.
Let $1\in\mathscr S_0$ act identically.  Reduction gives the identical
construction for $\mathcal M_0$.  Finiteness follows from the finite free
value complexes over the Noetherian ring $R$ and their finite-dimensional
reductions.
\end{proof}

The notation $V_{\mathrm s}^r$ refers to cohomology \emph{after} reducing the
complex.  It does not identify this group with $V^r/tV^r$; the possible
torsion contribution will be retained in the lifting argument.

\subsection{Analytic theta functions and prescribed field embeddings}

Consider the nonempty convex domain
\begin{equation}\label{theta:period-domain}
 \mathcal P_M=
 \left\{(\tau,s)\in\C^2:
    \Pi=\tau1_8+\frac{sM}{2\pi i},\quad\Im\Pi>0\right\}.
\end{equation}
Its periods are symmetric and commute with $D$.  Nonemptiness follows
by taking $s=0$ and $\Im\tau>0$; convexity is the convexity of positive
definite real symmetric matrices.  The following theta description applies
to every symmetric $\Pi$ with $\Im\Pi>0$, in particular to the whole
marked domain $\mathcal U$.  Denote by $\Theta_\Pi$ the theta line
on $A_\Pi=\C^8/(\Z^8+\Pi\Z^8)$ with multipliers
\[
 F(z+b)=F(z),\qquad
 F(z+\Pi b)=e^{-\pi i b^t\Pi b-2\pi i b^tz}F(z),
 \qquad b\in\Z^8.
\]
For $m>0$ a basis of its $m$th power is
\begin{equation}\label{theta:analytic-basis}
 \vartheta_{m,\bar u}(z,\Pi)
 =\sum_{u\in\bar u+\Z^8}
       \exp\bigl(\pi i m u^t\Pi u+2\pi i m u^tz\bigr),
 \qquad \bar u\in(m^{-1}\Z^8)/\Z^8.
\end{equation}
Positive definiteness gives normal Gaussian convergence on compact sets.
The multiplier equations follow by reindexing.  Conversely a periodic
holomorphic section has a Fourier expansion in $e^{2\pi i n^tz}$; its
second multiplier equation relates the coefficients whose indices differ
by $m\Z^8$.  One freely chosen coefficient for each coset gives exactly
Equation~\eqref{theta:analytic-basis}, proving that this is the full basis.
The convergence is normal jointly in $z$ and $\Pi$, so these sections
also vary holomorphically over $\mathcal U$.
Restrict now to $\mathcal P_M$.  Completing the square in the product gives
Equation~\eqref{theta:product} with substitutions
\begin{equation}\label{theta:analytic-substitution}
 \qnov^\lambda\rightsquigarrow e^{2\pi i\tau\lambda},
 \qquad t\rightsquigarrow s.
\end{equation}
Indeed the extra exponent is
$2\pi i\tau T^{1_8}_{m,m'}+sT^M_{m,m'}$.

These substitutions currently describe analytic functions, not a map
on the entire formal coefficient field.  The following statement supplies
exactly the field map that is needed.

\begin{proposition}[Realizing the theta constants]\label{theta:realization}
There is a countable union $E_M$ of proper closed analytic subsets of
$\mathcal P_M$ such that every $(\tau,s)\in\mathcal P_M\setminus E_M$
admits a field embedding
\[
 \iota_{\tau,s}:\operatorname{Frac}(R)\hookrightarrow\C
\]
under which
\begin{equation}\label{theta:generic-section-ring}
 \mathscr S\otimes_{R,\iota_{\tau,s}}\C
 \ \simeq\
 \bigoplus_{\ell\geq0}
       H^0(A_\Pi,\Theta_\Pi^{\,h_0\ell}).
\end{equation}
The isomorphism sends $e_{\ell,\bar u}$ to
$\vartheta_{h_0\ell,\bar u}$.  Likewise, for some
$\tau$ in the upper half-plane an embedding $\kappa\hookrightarrow\C$
identifies $\mathscr S_{\mathrm s}\otimes_\kappa\C$ with the same section
ring for $\Pi=\tau1_8$.
\end{proposition}

\begin{proof}
List the countably many scalar multiplication coefficients of
Equation~\eqref{theta:product} as $c_1,c_2,\ldots\in\Lambda[[t]]$,
and their analytic functions as $c_1^{\mathrm{an}},c_2^{\mathrm{an}},\ldots$.
We first prove that, for every polynomial $P\in\Q[X_1,\ldots,X_N]$,
\begin{equation}\label{theta:relations}
 P(c_1,\ldots,c_N)=0
 \quad\Longleftrightarrow\quad
 P(c_1^{\mathrm{an}},\ldots,c_N^{\mathrm{an}})\equiv0
       \text{ on }\mathcal P_M.
\end{equation}
Only finitely many integers $m,m'$ occur in this expression.
Their area exponents are nonnegative rational numbers with a common
denominator $D_0$.  This follows directly from
Equation~\eqref{theta:quadratic}, since the input lattice denominators
are $m$ and $m'$.  Multiplying finitely many such series preserves the
common denominator.  Put $Q=e^{2\pi i\tau/D_0}$.  There is a constant
$C_M$ such that
\[
 |T^M_{m,m'}(u,v)|\leq C_M T^{1_8}_{m,m'}(u,v)
\]
for all these summands; the operator norm bound for $M$ even gives a
constant independent of $m,m'$.  Thus the series and all their finite
products converge normally on a sufficiently small bidisk in $(Q,s)$
centered at $(0,0)$.  For example their absolute terms are bounded by
$|Q|^{D_0T^{1_8}}e^{|s|C_MT^{1_8}}$, with only polynomially many lattice
points of bounded area.  They extend holomorphically across $Q=0$.

The resulting joint power series in $Q,s$ is exactly the series obtained
formally by expanding $\exp(tT^M)$ and replacing
$\qnov^{n/D_0}$ by $Q^n$.  For any coefficient of $t^a$, distinct
$\qnov^{n/D_0}$ have uniquely determined coefficients in $\Lambda$.
The injection $\Lambda\hookrightarrow\kappa$ preserves that equality.
Consequently vanishing in $\kappa[[t]]$ is equivalent to vanishing of
this joint holomorphic germ.  On the punctured bidisk it is the analytic
theta expression, with a branch of $\tau$ of sufficiently large imaginary
part.  The identity theorem on connected $\mathcal P_M$ proves
Equation~\eqref{theta:relations}.

Let $F=\Q(c_1,c_2,\ldots)\subset\operatorname{Frac}(R)$.  It is a
countable field.  For each rational polynomial whose formal value is
nonzero, exclude the zero set of its analytic value.  By
Equation~\eqref{theta:relations} these are countably many proper closed
analytic subsets.  At every remaining point, evaluation defines an
injective homomorphism on $\Q[c_1,c_2,\ldots]$ and hence an embedding
$F\hookrightarrow\C$; all nonzero denominators remain nonzero.

To extend this embedding, choose a transcendence basis of
$\operatorname{Frac}(R)$ over $F$.  Its cardinality is at most
$|\C|$, since
$|\kappa((t))|\leq|\C|^{\aleph_0}=|\C|$.
Over the countable image of $F$, the field $\C$ has transcendence degree
$|\C|$: the algebraic closure of a field generated by an infinite set of
cardinality $\lambda$ over a countable field has cardinality
$\max(\lambda,\aleph_0)$.  Choose that many algebraically independent
target elements, map the basis to them, and extend across the remaining
algebraic extension using algebraic closedness of $\C$.  This gives
$\iota_{\tau,s}$.  Its prescribed values on all multiplication constants,
together with the theta bases, give
Equation~\eqref{theta:generic-section-ring} degree by degree.

For the special fiber, perform the same argument with the constants
$c_i|_{t=0}$ and the one-variable cusp expansion in $Q$.
Their analytic functions are the structure constants at $s=0$.
Avoid their countably many nonidentical relations in the upper half-plane
and extend the resulting countable-field embedding to $\kappa$ by the
same cardinality argument.
\end{proof}

\begin{remark}\label{theta:abstract-embedding}
The embedding in Proposition~\ref{theta:realization} prescribes the values
of whole theta structure constants.  It need not fix the original complex
coefficient field, send the formal element $t$ to the numerical parameter
$s$, or preserve the Novikov valuation.  None of those additional
requirements is used.  In particular the proposition makes no analytic
evaluation of the new brane's Floer series.  It remains valid after any
enlargement of $\kappa$ that retains the injection of $\Lambda$ and the
cardinality bound.  The same constants and the same excluded analytic
sets can be used.
\end{remark}

\subsection{Normal generation and the family over the trait}

We include the needed normal-generation argument to specify the algebraic
properties of the special fiber without making a finite-generation
assumption on the formal ring.

\begin{lemma}[Elliptic normal generation]\label{theta:elliptic-generation}
Let $E$ be a smooth complex elliptic curve and $L$ a line bundle of degree
$h\geq3$.  Then $L$ is very ample and every multiplication map
\[
 H^0(E,L)\otimes H^0(E,L^a)\longrightarrow H^0(E,L^{a+1}),
 \qquad a\geq1,
\]
is surjective.
\end{lemma}

\begin{proof}
Riemann--Roch and Serre duality give $H^1(E,A)=0$ and
$h^0(E,A)=\deg A$ for every positive-degree line bundle $A$.
For every subscheme $Z\subset E$ of length two,
$\deg L(-Z)>0$, so evaluation to $L|_Z$ is surjective.
This separates both points and tangent vectors and proves very ampleness.

First consider $a=1$.  If $h=3$, the complete linear system embeds $E$
as a plane cubic.  No nonzero quadratic polynomial vanishes on that
integral cubic, so
$\operatorname{Sym}^2H^0(E,L)\to H^0(E,L^2)$ is injective.
Both sides have dimension six, and the map is therefore surjective.
For $h>3$, apply induction to $L(-p)$ for any point $p$.
Products of its sections span $H^0(E,L^2(-2p))$.
There are sections $s_0,s_1$ of $L$ with vanishing orders zero and one at
$p$, respectively, by the same positive-degree cohomology calculation.
Modulo sections vanishing to order two, the products $s_0^2,s_0s_1$
span the two-dimensional space of first jets of $L^2$ at $p$.
They complete the image to $H^0(E,L^2)$.

Choose a basepoint-free pencil $W\subset H^0(E,L)$: take a nonzero
section and then a second section nonvanishing at each of its finitely
many zeros.  The pencil sequence is
\[
 0\longrightarrow L^{-1}\longrightarrow W\otimes\cO_E
       \longrightarrow L\longrightarrow0.
\]
After tensoring by $L^a$, with $a\geq2$, its first term is $L^{a-1}$
and has zero $H^1$.  Hence even multiplication by $W$ surjects onto
$H^0(E,L^{a+1})$.
\end{proof}

At $\Pi=\tau1_8$, the variety $A_\Pi$ is the product of eight elliptic
curves and $\Theta_\Pi^{h_0}$ is the external tensor product of degree
$h_0$ lines.  K\"unneth's formula for sections and
Lemma~\ref{theta:elliptic-generation} give generation of its section ring
in degree one.  The product embedding followed by the Segre embedding
is the embedding of its complete linear system.  Its section-ring Proj
is therefore that smooth integral eightfold, with grading line
$\Theta_\Pi^{h_0}$.

\begin{proposition}[The smooth theta family]\label{theta:family}
The ring $\mathscr S$ is a Noetherian graded $R$-algebra generated by the
finite free module $\mathscr S_1$.  Its relative Proj
\[
 \mathscr X=\Proj\mathscr S\longrightarrow\Spec R
\]
is smooth, projective, and flat, and its total space is regular.
The special fiber
$\mathscr X_{\mathrm s}=\Proj\mathscr S_{\mathrm s}$ is a smooth
geometrically integral eightfold.  Its grading line $\cO(1)$ is very
ample, and the canonical map
\[
 \mathscr S_{\mathrm s}\longrightarrow
 \bigoplus_{\ell\geq0}H^0(\mathscr X_{\mathrm s},\cO(\ell))
\]
is an isomorphism.  Under any realization in
Proposition~\ref{theta:realization}, the geometric generic fiber and its
grading line become $(A_\Pi,\Theta_\Pi^{h_0})$.
\end{proposition}

\begin{proof}
Apply the special-fiber embedding of
Proposition~\ref{theta:realization}.  After the faithfully flat field
extension $\kappa\hookrightarrow\C$, the preceding elliptic argument
identifies the ring, Proj, and grading line with the claimed product
and its full section ring.  Surjectivity of multiplication descends
degree by degree.  Formation of Proj commutes with field extension;
smoothness, geometric integrality, and the complete-linear-system
embedding descend as well.  Proper coherent cohomology commutes with
field extension, so the canonical section-ring isomorphism descends.

For each $\ell$, the cokernel of
$\operatorname{Sym}^\ell_R\mathscr S_1\to\mathscr S_\ell$ is a finite
$R$-module.  Its reduction modulo $t$ is zero by special-fiber normal
generation.  Nakayama's lemma makes the cokernel zero.  Thus $\mathscr S$
is a quotient of a polynomial ring on the finite free module
$\mathscr S_1$ over the Noetherian DVR $R$.  It is Noetherian and its
Proj is projective, with $\cO(1)$ relatively very ample.

The direct sum of the free graded pieces makes $\mathscr S$ flat over
$R$.  For a homogeneous element $f$, the affine chart ring
$(\mathscr S[f^{-1}])_0$ is the degree-zero direct summand of an
$R$-flat localization.  It is flat; these charts prove flatness of
$\mathscr X$.  A flat morphism of finite presentation with geometrically
smooth special fiber is smooth along that fiber.  The complement of
the smooth locus is closed, and its image is closed because $\mathscr X$
is proper over $\Spec R$.  If it were nonempty, this image would contain
the closed point of the trait, contradicting the preceding smoothness.
Hence the morphism is smooth everywhere.  The base is regular, so its
smooth total space is regular.  Finally the generic-fiber assertion
follows by taking Proj in Equation~\eqref{theta:generic-section-ring}.
\end{proof}
